\bookchapter{Kapitel 29: Digitale Gesundheit – Eine gesunde Beziehung zum Smartphone aufbauen}

\chapterintro{Kennen Sie das: Sie wollten nur 5 Minuten Kurzvideos schauen, schauen auf – und zwei Stunden sind vergangen? Oder Sie liegen abends im Bett, scrollen durchs Handy, werden immer wacher und schlafen erst um zwei Uhr nachts ein? Das Handy hat unser Leben verändert, aber es verschlingt unauffällig unsere Zeit und Gesundheit. Dieses Kapitel will Sie nicht dazu bringen, “das Handy aufzugeben”, sondern Ihnen helfen, die Kontrolle zurückzugewinnen.}

\booksection{I. Schauen Sie zuerst auf Ihre “Bildschirmzeit” —— die Zahlen öffnen die Augen}

Bevor Sie etwas ändern, müssen Sie den Ist-Zustand kennen. Das Handy hat eine Funktion, die viele nie geöffnet haben: “Bildschirmzeit” (iPhone) oder “Digitale Gesundheit” (Android). Sie zeigt, wie viel Zeit Sie täglich auf dem Handy verbringen und wofür.

\begin{itemize}
  \item \textbf{iPhone}: Einstellungen > Bildschirmzeit > Alle Aktivitäten anzeigen. Sie sehen tägliche Entsperrungen, wie oft Sie das Handy hochnehmen und die Nutzungsdauer jeder App.
  \item \textbf{Android}: Einstellungen > Digitale Gesundheit und Kindersicherung. Bei Huawei "Gesunde Handy-Nutzung", bei Xiaomi "Bildschirmzeit-Verwaltung". Ähnliche Funktionen, alle zeigen tägliche Nutzungsdaten.
\end{itemize}

\textbf{80 \% derer, die diese Daten zum ersten Mal sehen, sind schockiert}. Ein Freund stellte fest, dass er täglich über 4 Stunden auf Douyin (TikTok) verbringt – das sind 28 Stunden pro Woche, so viel wie ein Nebenjob. Dieser “digitale Schock” ist der erste Schritt zur Veränderung.

\booksection{II. App-“Nutzungslimits” setzen —— einen Geschwindigkeitsbegrenzer einbauen}

Wenn Sie wissen, wo das Problem liegt, können Sie gezielt Limits setzen. Sie müssen die App nicht löschen – geben Sie ihr einfach ein “tägliches Kontingent”:

\begin{enumerate}
  \item \textbf{iPhone}: Einstellungen > Bildschirmzeit > App-Limits > Limit hinzufügen. Wählen Sie die App (z. B. Douyin), setzen Sie täglich 30 Minuten. Danach wird das App-Symbol grau, und zum Öffnen ist ein zusätzliches Passwort nötig.
  \item \textbf{Android}: Einstellungen > Digitale Gesundheit > App-Timer. App auswählen und tägliches Limit setzen – danach wird das Symbol grau.
  \item \textbf{Tipp}: Setzen Sie das Limit-Passwort auf eines, das Sie selten verwenden und schwer zu merken ist. Ist es zu leicht, werden Sie jedes Mal "verlängern" – und das Limit ist wirkungslos.
\end{enumerate}

\booksection{III. “Nicht stören” richtig nutzen —— nicht Benachrichtigungen ausschalten, sondern Stress}

Viele Handys piepen dutzende Male am Tag, und jedes Mal wird die Aufmerksamkeit unterbrochen. Sie denken, Sie “schauen nur kurz” – aber nach jeder Unterbrechung braucht man 15–20 Minuten, um sich wieder zu konzentrieren. “Nicht stören” heißt nicht “unerreichbar sein”, sondern hilft Ihnen, die Aufmerksamkeit zurückzugewinnen.

\begin{itemize}
  \item \textbf{Zeitgesteuert aktivieren}: "Nicht stören" von 22:00 bis 7:00 Uhr automatisch aktivieren. So weckt Sie keine Spam-SMS im Schlaf.
  \item \textbf{Wichtige Kontakte durchlassen}: In den "Nicht stören"-Einstellungen "Anrufe erlauben" auf "Favoriten" oder "Wiederholte Anrufe" setzen (zweiter Anruf innerhalb von 3 Minuten klingelt normal). So verpassen Sie keine Notrufe der Familie, werden aber nicht von Werbeanrufen gestört.
  \item \textbf{Arbeits-Fokus-Modus}: Der "Fokus" des iPhone lässt benutzerdefinierte Szenen einstellen – z. B. im Modus "Arbeit" nur arbeitsbezogene App-Benachrichtigungen, im Modus "Lesen" alle Benachrichtigungen stumm.
\end{itemize}

\booksection{IV. Augenschutz durch Gewohnheit, nicht durch Apps —— die 20-20-20-Regel}

Wer lange aufs Handy schaut, bekommt trockene und schmerzende Augen – das ist unvermeidlich. Der Augenschutzmodus (Blaufilter) hilft ein wenig, aber wirklich wirksam ist, die Sehgewohnheiten zu ändern:

\begin{tcbinfo}{Die 20-20-20-Regel}
Alle\textbf{20 Minuten}auf Bildschirm schauen, dann den Blick auf\textbf{20 Fuß (ca. 6 m)}in die Ferne richten und\textbf{20 Sekunden}halten. Diese einfache Übung lindert die Ermüdung des Ziliarmuskels wirksam.

\end{tcbinfo}

\begin{itemize}
  \item \textbf{Augenschutzmodus aktivieren}: Einstellungen > Anzeige > Augenschutzmodus (oder Nachtmodus). Aktiviert wird der Bildschirm wärmer und reduziert Blaues Licht. Empfehlung: Automatisch von Sonnenuntergang bis Sonnenaufgang aktivieren.
  \item \textbf{Bildschirmhelligkeit anpassen}: Ist es dunkel im Raum, sollte auch der Bildschirm dunkel sein. Viele haben in Dunkelheit die Helligkeit zu hoch – das ist das schädlichste Verhalten für die Augen. "Auto-Helligkeit" aktivieren und kontrollieren, dass sie funktioniert.
  \item \textbf{Abstand halten}: Mindestens 30 cm Abstand zwischen Handy und Augen. Viele halten das Handy zu nah – der Ziliarmuskel ist ständig angespannt, und die Kurzsichtigkeit verschlimmert sich schnell.
\end{itemize}

\booksection{V. Nacken-Rettung für die “Kopf-nach-unten-Fraktion”}

Wenn Sie mit gesenktem Kopf aufs Handy schauen, lastet das 5–6-fache des normalen Drucks auf Ihrem Nacken. Langfristig führt das zu einer Begradigung der Halswirbelsäule, Knochenausbildungen, Schwindel und tauben Händen. Hier sind einfache Verbesserungsmaßnahmen:

\begin{itemize}
  \item \textbf{Handy höher halten}: Halten Sie das Handy beim Schauen möglichst auf Augenhöhe, statt mit gesenktem Kopf auf das Handy auf den Knien zu blicken. Der Arm wird müde, aber der Nacken dankt es.
  \item \textbf{Alle 30 Minuten den Nacken bewegen}: Langsam den Kopf nach links drehen bis zum Anschlag, 3 Sekunden halten, zurück; dann nach rechts, 3 Sekunden halten. 5-mal wiederholen. Danach das "Reiszeichen-Mal-Exercise" – mit dem Kinn das Zeichen „Mi“ in die Luft schreiben.
  \item \textbf{3 Minuten an der Wand stehen}: Hinterkopf, Schulterblätter, Gesäß und Fersen an die Wand pressen – täglich 3 Minuten stehen korrigiert wirksam die Haltung.
\end{itemize}

\booksection{VI. Schlaf und Handy —— die Gewohnheit, die Sie im Schlafzimmer am meisten ändern sollten}

Das blaue Licht des Handybildschirms hemmt die Melatoninausschüttung und macht Sie “immer wacher”. Viele meinen: “Ich werde beim Scrollen müde”, aber in Wirklichkeit täuscht das blaue Licht das Gehirn – wenn Sie das Handy weglegen, können Sie nicht einschlafen. Verbesserungsvorschläge:

\begin{itemize}
  \item \textbf{1 Stunde vor dem Schlafengehen nicht aufs Handy schauen}: Die wirksamste und schwerste Regel. Wenn es nicht klappt, zumindest den "Nachtmodus/Augenschutzmodus" aktivieren, um Blaues Licht zu minimieren.
  \item \textbf{Das Handy außer Reichweite legen}: Ladegerät am anderen Ende des Schlafzimmers platzieren, nicht auf dem Nachttisch. So greifen Sie nicht unbändig zum Handy und müssen aufstehen, um den Wecker auszuschalten – zwei Fliegen mit einer Klappe.
  \item \textbf{"Schlafmodus" aktivieren}: Der "Schlaf"-Fokus des iPhone kann den Bildschirm während der Schlafzeit automatisch sperren und den Sperrbildschirm vereinfachen; Android-Marken bieten ähnliche Funktionen.
  \item \textbf{Echten Wecker statt Handy verwenden}: Kaufen Sie für ein paar Yuan einen echten Wecker. Nutzen Sie nicht das Handy als Wecker – Sie entgehen der Falle "nur kurz die Uhr checken und dann doch scrollen".
\end{itemize}

\booksection{VII. Informationsangst —— “selektive Blindheit” lernen}

Wir nehmen täglich so viele Informationen auf wie ein Mensch früher in seinem ganzen Leben. Aber das Gehirn ist nicht dafür gemacht, so viel zu verarbeiten – Informationsüberflutung führt zu Angst, Konzentrationsverlust und Entscheidungsmüdigkeit. Die Lösung ist nicht “mehr Informationen”, sondern “aktives Ausblenden”:

\begin{itemize}
  \item \textbf{Abonnements kündigen}: Öffnen Sie die Liste aller von Ihnen abonnierten öffentlichen Accounts, UP-Master und Blogger. Fragen Sie sich: "Habe ich im letzten Monat tatsächlich einen Wert daraus gezogen?" Wenn nein, abbestellen. Qualität vor Quantität.
  \item \textbf{Nicht benötigte App-Pushs deaktivieren}: Außer WeChat, Telefon, SMS und Kalender alle anderen App-Benachrichtigungen ausschalten. Nachrichten brauchen keinen Push – Sie sehen sie, wenn Sie aktiv danach suchen.
  \item \textbf{Einen "Informationsfastentag" einführen}: Wählen Sie einen Tag pro Woche (z. B. Sonntag), an dem Sie keine Handy-News und keine Kurzvideos konsumieren. Sie werden merken: Die Welt geht nicht unter, und am Montag sind Sie deutlich fitter.
\end{itemize}

\begin{tcbpractice}{Praxis: Eine Woche digitale Gesundheit}
\begin{enumerate}
  \item Heute: Bildschirmzeit/Digitale Gesundheit öffnen und die Daten der letzten Woche ansehen.
  \item Morgen: Für die zeitintensivste Unterhaltungs-App ein tägliches Limit setzen (empfohlen 30–60 Minuten).
  \item Übermorgen: "Nicht stören" zeitgesteuert aktivieren (22:00–07:00), Favoriten-Anrufe erlauben.
  \item Tag 4: 1 Stunde vor dem Schlafengehen nicht aufs Handy schauen, Ladegerät ans andere Ende des Schlafzimmers legen.
  \item Tag 5: 5 öffentliche Accounts oder Blogger abbestellen, die Sie nicht mehr lesen.
  \item Tag 6: Einen ganzen Tag "Informationsfasten" – keine News, keine Kurzvideos.
  \item Tag 7: Den Bildschirmzeit-Bericht dieser Woche ansehen und mit der vorherigen vergleichen.
\end{enumerate}

\end{tcbpractice}

\begin{tcbsummary}{Kapitelzusammenfassung}
\begin{itemize}
  \item Zuerst die Bildschirmzeit-Daten ansehen – den Ist-Zustand erkennen ist der erste Schritt
  \item Unterhaltungs-Apps Limits setzen, "Nicht stören" zeitsteuern – die Aufmerksamkeit zurückgewinnen
  \item Die 20-20-20-Regel ist die einfachste und wirksamste Augengewohnheit
  \item Bei gesenktem Kopf lastet das 5–6-fache auf dem Nacken – Handy höher halten + regelmäßig den Nacken bewegen
  \item 1 Stunde vor dem Schlafengehen nicht aufs Handy schauen, Ladegerät außer Reichweite legen
  \item Das Gegenmittel gegen Informationsangst ist "aktives Ausblenden", nicht "mehr konsumieren"
\end{itemize}

\end{tcbsummary}
